\documentclass[letterpaper,10pt]{article}

% Editable source for Pooja_Kiran_Security_Engineer_Resume.pdf.
% Numbers in the projects/summary sections are reconciled to the verified
% VERIFIED_METRICS.md evidence in the three flagship repositories
% (MCP 648, AWS Identity Guard 235, HF scanner 211; total 1,094).

\usepackage[margin=0.5in]{geometry}
\usepackage{titlesec}
\usepackage{enumitem}
\usepackage[hidelinks]{hyperref}
\usepackage{xcolor}
\usepackage{parskip}

\definecolor{rulegray}{gray}{0.55}
\setlength{\parindent}{0pt}
\pagestyle{empty}

\titleformat{\section}{\large\bfseries}{}{0pt}{}[\vspace{-0.6em}\color{rulegray}\rule{\linewidth}{0.6pt}]
\titlespacing{\section}{0pt}{0.5em}{0.25em}

\newcommand{\entry}[2]{\textbf{#1}\hfill #2\par}
\newcommand{\subentry}[1]{\textit{#1}\par}
\setlist[itemize]{leftmargin=1.4em,topsep=1pt,itemsep=1pt,parsep=0pt}

\begin{document}

\begin{center}
  {\huge\bfseries POOJA KIRAN}\\[2pt]
  {\large\bfseries SECURITY ENGINEER}\\[3pt]
  Tempe, AZ | +1 480-776-7745 | poojakiranbhardwaj@gmail.com | linkedin.com/in/poojakiran | github.com/poojakira
\end{center}

\section*{PROFESSIONAL SUMMARY}
Security Engineer with hands-on experience owning security architecture, control development, adversarial
validation, and security automation across AI agents, AWS IAM, and AI model supply chains. Architected three
open-source Python security systems that address unsafe tool execution, excessive agent privileges, and model
artifacts from unverified sources, backed by 1{,}094 documented passing tests across the three repositories.
Translates threat scenarios into enforceable controls, CI gates, and detection workflows using MCP, AWS IAM,
SARIF, GitHub Code Scanning, Elastic Security, and SIEM telemetry.

\section*{TECHNICAL SKILLS}
\textbf{Security Engineering \& Testing:} Security Engineering, Information Security, Security Solutions, Security Automation, Security Testing, Automated Testing, Threat Modeling, Vulnerability Scanning, Incident Response Runbooks\\
\textbf{Programming \& Platforms:} Python, Bash/Shell, FastAPI, REST APIs, JSON-RPC 2.0, Linux, Docker\\
\textbf{AI \& Agent Security:} Model Context Protocol (MCP), Prompt Injection Detection, Secure Tool Invocation, Tool-Call Inspection, Capability Enforcement, PII Leakage Detection, Data Exfiltration Detection\\
\textbf{Cloud \& Identity Security:} AWS IAM, Least Privilege, IAM Policy Analysis, iam:PassRole, sts:AssumeRole, Trust Policies, Permission Boundaries, Privilege Escalation Analysis, CloudTrail\\
\textbf{Network \& Detection Security:} Network Security Controls, Network Egress Controls, Elastic Security, SIEM, Detection Engineering, Security Telemetry, MITRE ATT\&CK, Audit Logging\\
\textbf{DevSecOps \& AI Supply Chain:} GitHub Actions, CI/CD Security, SARIF 2.1.0, CodeQL, Bandit, Trivy, pip-audit, Model Provenance, Model Artifact Security, Pickle Analysis, SafeTensors, GGUF, ONNX, SBOM

\section*{WORK EXPERIENCE}
\entry{Independent AI Security Researcher \& Engineer}{Jul 2024 - Present}
\subentry{Self-Directed Security Research | Tempe, AZ}
\begin{itemize}
  \item Owned the architecture and delivery of three open-source Python security systems addressing unsafe AI tool invocation, excessive agent privileges, and model artifacts from unverified sources, codifying 55 prompt-injection patterns and 25 IAM security rules.
  \item Drove end-to-end validation through 1{,}094 documented passing tests across the three repositories, establishing measurable assurance for prompt-injection controls, IAM rule behavior, model-artifact inspection, parser edge cases, and adversarial scenarios.
  \item Integrated security findings into automated engineering workflows through SARIF 2.1.0, GitHub Code Scanning, Elastic Security detections, SIEM validation, and CI-integrated enforcement gates.
\end{itemize}

\entry{Business and Compliance Lead | AEROSEC Externship}{Jul 2025 - Nov 2025}
\subentry{Arizona State University | Honeywell Aerospace Partnership | Tempe, AZ}
\begin{itemize}
  \item Directed business, security, and compliance planning for a \$120K first-year passenger-service-system (PSS) security scenario, translating third-party cybersecurity risk into deployment, operational readiness, and commercialization decisions.
  \item Owned a five-year financial and deployment model evaluating implementation cost, growth assumptions, and security trade-offs, identifying a scenario with approximately 12\% cost savings.
  \item Presented security, compliance, financial, and deployment recommendations to ASU faculty and Honeywell mentors, supporting risk-informed project strategy and operational readiness decisions.
\end{itemize}

\entry{Graduate Teaching Assistant | IT Grader}{Dec 2024 - Sep 2025}
\subentry{Arizona State University | Tempe, AZ}
\begin{itemize}
  \item Managed technical assessment of approximately 85 undergraduate web programming submissions per term, identifying implementation defects and delivering structured, rubric-driven engineering feedback.
  \item Evaluated 52 graduate IT security assignments spanning cybersecurity policy, compliance, and security concepts while enforcing consistent technical assessment standards.
  \item Resolved grading inconsistencies with ASU faculty and communicated technical corrections across programming and cybersecurity coursework.
\end{itemize}

\section*{SELECTED SECURITY ENGINEERING PROJECTS}
\entry{MCP Agent Security Gateway | \href{https://github.com/poojakira/mcp-agent-security-gateway}{GitHub}}{Jul 2026 - Present}
\begin{itemize}
  \item Closed an agent-security gap by architecting an inline MCP and JSON-RPC 2.0 gateway that inspects tool calls before downstream execution, creating a policy enforcement boundary between agent intent and privileged actions.
  \item Hardened the execution boundary with 55 prompt-injection patterns, Unicode/Base64/ROT13 normalization, capability checks, PII and exfiltration signals, hash-chained audit logging, rate limiting, and security telemetry.
  \item Validated 648 passing tests at 79.6\% statement coverage and defined 9 MITRE ATT\&CK-mapped Elastic Security rules, 21 SIEM tests, and end-to-end ELK event shipping with 0 Filebeat errors in the local Docker lab.
\end{itemize}

\entry{AWS Agent Identity Guard | \href{https://github.com/poojakira/aws-agent-identity-guard}{GitHub}}{Aug 2026 - Sep 2026}
\begin{itemize}
  \item Reduced pre-deployment IAM risk for autonomous agents by architecting a static analyzer that identifies excessive and high-risk permissions, privilege-escalation paths, risky trust policies, audit tampering, and missing permission boundaries before role deployment.
  \item Codified 25 deterministic rules and converted findings into CI-enforceable JSON and SARIF 2.1.0 output, GitHub Code Scanning annotations, and exit-code gates for high and critical conditions.
  \item Validated all 25 rule IDs across 235 passing tests, including positive and negative cases, parser fuzzing, SARIF conformance, and failure-mode tests.
\end{itemize}

\entry{HF Model Provenance Scanner | \href{https://github.com/poojakira/hf-model-provenance-scanner}{GitHub}}{Jul 2026 - Sep 2026}
\begin{itemize}
  \item Addressed AI model supply-chain risk by engineering a non-executing scanner that inspects untrusted repositories and artifacts for provenance gaps, unsafe serialization, suspicious loader behavior, impersonation, and configuration anomalies before model loading.
  \item Expanded inspection across Python loaders, pickle bytecode, SafeTensors, GGUF, ONNX, dependency indicators, signatures, and SBOM evidence without executing untrusted artifacts.
  \item Validated 211 passing tests; committed red-team fixture suites detected 12/12 core incident reproductions and 18/18 extended variants with 0 actionable false positives across 4 benign samples.
\end{itemize}

\section*{EDUCATION}
\entry{Master of Science in Information Technology}{Aug 2024 - May 2026}
\subentry{Arizona State University | Tempe, AZ | GPA: 3.87/4.00}
\entry{Bachelor of Engineering in Computer Science and Engineering}{Aug 2019 - Aug 2023}
\subentry{Ramaiah University of Applied Sciences | Bengaluru, KA, India}

\section*{TRAINING \& CREDENTIALS}
\textbf{AWS Academy Graduate - Cloud Security Foundations} | AWS Academy\\
\textbf{AWS Academy Graduate - Cloud Architecting} | AWS Academy

\section*{PUBLICATIONS}
\textbf{A Personalized E-Learning System Using Reinforcement Learning Through Satellite} | IEEE INDICON 2023 | Dec 2023\\
Co-authored and presented a peer-reviewed IEEE conference paper on reinforcement-learning-based personalized e-learning for remote communities using satellite connectivity.

\section*{AWARDS \& GRANTS}
\entry{KSCST 46th Series Student Project Programme Research Grant}{Apr 2023}
\subentry{Karnataka State Council for Science and Technology (KSCST)}
Awarded INR 6{,}000 in project funding for an AI-based e-learning initiative for remote communities.

\end{document}
